\documentclass{report}
\usepackage{amsmath,amssymb}
\setlength{\parindent}{0mm}
\setlength{\parskip}{1em}
\begin{document}
\begin{center}
\rule{6in}{1pt} \
{\large ILYA LASHUK \\
{\bf COARSE DE RHAM COMPLEXES WITH IMPROVED APPROXIMATION PROPERTIES}}

7000 East Ave \\ L-560 \\ LIVERMORE \\ CA \\ 94550
\\
{\tt LASHUK2@LLNL.GOV}\\
PANAYOT VASSILEVSKI\end{center}

\newcommand{\dvg}{\mathbf{div}}
\newcommand{\curl}{\mathbf{curl}}
\def\Nedelec{N\'ed\'elec}

We present a technique based on element agglomeration for constructing
coarse subspaces of:
(a) the lowest order tetrahedral \Nedelec~ space, and (b) the scalar
piecewise-linear finite element
space, in both cases assuming general unstructured mesh. The constructed
coarse spaces, together with coarse $H(\dvg)$ and $L_2$ spaces described
in our previous work,
form an exact sequence with respect to the exterior
derivative operators (for domains homeomorphic to a ball). The constructed
coarse counterparts of the \Nedelec~ and Raviart--Thomas spaces locally
contain the vector constant functions,
whereas the respective coarse version of the scalar piecewise-linear
space locally contains all linear functions.
These properties hold regardless of the shape of the agglomerates, as
long as each agglomerate stays
homeomorphic to a ball (which is easily ensured in practice).
We illustrate the approximation properties of the constructed coarse
spaces as well as the convergence
of respective two-level AMGe methods associated with them.


\end{document}
